\subsection{Observable guidance with energies, electron counts, forces, and stress}
\label{subsec:training_objectives}

\texttt{Mandala} supports pure matrix learning, mixed matrix-and-observable learning, and
fully observable-guided training from one common operator-prediction pipeline.
For a chosen set of matrix targets $\mathcal{M} \subseteq \{H,S,D\}$, the
general training objective is
\begin{equation}
\Loss
=
\sum_{X \in \mathcal{M}} \lambda_X \Loss_X^{\mathrm{matrix}}
+ \lambda_E \Loss^{E}
+ \lambda_N \Loss^{N}
+ \lambda_F \Loss^{F}
+ \lambda_{\sigma} \Loss^{\sigma}.
\end{equation}
The matrix terms compare predicted and reference sparse operators, typically by
a weighted combination of $\ell_1$ and $\ell_2$ errors,
\begin{equation}
\Loss_X^{\mathrm{matrix}}
=
\eta \, \mathrm{MAE}(\widehat{X},X)
+ (1-\eta)\, \mathrm{MSE}(\widehat{X},X),
\end{equation}
with $\eta$ controlled by configuration.

Observable guidance is defined on quantities derived from the predicted
operators. In the operator-centered workflow used throughout \texttt{Mandala}, the
energy observable used in training and the number of electrons are
\begin{equation}
\widehat{E} = \Tr(\widehat{D}\widehat{H}),
\qquad
\widehat{N}_e = \Tr(\widehat{D}\widehat{S}),
\end{equation}
which gives
\begin{equation}
\Loss^E = \mathrm{MSE}(\widehat{E},E^{\mathrm{ref}}),
\qquad
\Loss^N = \mathrm{MSE}(\widehat{N}_e,N_e^{\mathrm{ref}}).
\end{equation}
The framework also supports mixed observable guidance paths in which one of the
operators entering the trace is replaced by the reference operator, for example
$\Tr(\widehat{H} D^{\mathrm{ref}})$,
$\Tr(\widehat{D} H^{\mathrm{ref}})$,
$\Tr(\widehat{D} S^{\mathrm{ref}})$, or
$\Tr(D^{\mathrm{ref}} \widehat{S})$, enabling diagnostic and training modes
that separate errors due to different predicted operators.

Forces are obtained by automatic differentiation of the predicted scalar energy
with respect to atomic coordinates,
\begin{equation}
\widehat{\bm{F}}_I
=
-\frac{\partial \widehat{E}}{\partial \bm{R}_I},
\end{equation}
and the corresponding force loss is
\begin{equation}
\Loss^F
=
\frac{1}{3N_{\mathrm{atoms}}}
\sum_I
\left\|
\widehat{\bm{F}}_I - \bm{F}_I^{\mathrm{ref}}
\right\|_2^2.
\end{equation}
No separate force head is required: forces are constrained through the same
predicted operators that define the energy.

For periodic systems, stress guidance is obtained by differentiating the
predicted energy with respect to the cell matrix $h$. If
$\Omega = \det(h)$ is the cell volume, \texttt{Mandala} uses
\begin{equation}
\widehat{\sigma}
=
\frac{1}{\Omega} h^{\top}
\frac{\partial \widehat{E}}{\partial h},
\end{equation}
with optional symmetrization
\begin{equation}
\widehat{\sigma}
\leftarrow
\frac{1}{2}\left(\widehat{\sigma}+\widehat{\sigma}^{\top}\right),
\end{equation}
and defines the stress loss as
\begin{equation}
\Loss^{\sigma}
=
\frac{1}{9}
\left\|
\widehat{\sigma} - \sigma^{\mathrm{ref}}
\right\|_F^2.
\end{equation}
This completes the observable-guidance picture: one can train a single
operator-valued model under simultaneous constraints from matrix entries,
energies, electron counts, forces, and stress tensors.
